% baselines and/or neural/Transformer
For the text classification task, we implemented two models: a Convolutional Neural Network (CNN) and a Long Short-Term Memory (LSTM) network. Both models were trained on the AG News dataset using the same train/dev/test splits, preprocessing, and feature extraction/tokenization (except for inherent CNN feature extraction layers) steps, which are described in the previous section. Tokenized text then was fed to the following models for training and evaluation.

\subsection{CNN Model}
The CNN model consists of an embedding layer, followed by three parallel convolutional layers with n-gram window sizes of 3, 5, and 7, each having 100 filters. The output of the convolutional layers is then passed through ReLU activation function, followed by max-pooling layers to capture the most prominent features. Finally, the pooled features are concatenated and fed into a fully connected layer with a default dropout rate of 0.5 for classification into the four categories. An illustration of the CNN architecture is shown in Figure \ref{fig:cnn_architecture}.

% TODO: Add CNN architecture figure here

% \begin{figure}[ht]
%     \centering
%     \includegraphics[width=0.8\textwidth]{cnn_architecture.png}
%     \caption{Architecture of the CNN model for text classification.}
%     \label{fig:cnn_architecture}
% \end{figure}


\subsection{LSTM Model}
The LSTM  model consists of 

% class LSTM(nn.Module):
%     """LSTM model for text classification.

%     Supports both unidirectional and bidirectional LSTM with proper handling
%     of variable-length sequences using pack_padded_sequence.
%     """

%     def __init__(
%         self,
%         vocab_size: int,
%         embed_dim: int,
%         hidden_dim: int,
%         num_classes: int,
%         num_layers: int = 2,
%         bidirectional: bool = False,
%         dropout: float = 0.5,
%     ) -> None:
%         """Initialize LSTM model.

%         Args:
%             vocab_size: Size of the vocabulary.
%             embed_dim: Dimension of word embeddings.
%             hidden_dim: LSTM hidden dimension.
%             num_classes: Number of output classes.
%             num_layers: Number of stacked LSTM layers (default: 2).
%             bidirectional: Whether to use bidirectional LSTM (default: False).
%             dropout: Dropout probability applied to hidden states (default: 0.5).
%         """
%         super().__init__()
%         self.embedding = nn.Embedding(vocab_size, embed_dim, padding_idx=0)
%         self.lstm = nn.LSTM(
%             embed_dim,
%             hidden_dim,
%             num_layers=num_layers,
%             batch_first=True,
%             bidirectional=bidirectional,
%             dropout=dropout if num_layers > 1 else 0,
%         )
%         self.dropout = nn.Dropout(dropout)
%         self.fc = nn.Linear(hidden_dim * (2 if bidirectional else 1), num_classes)

%     def forward(
%         self, x: torch.Tensor, lengths: torch.Tensor | None = None
%     ) -> torch.Tensor:
%         """Forward pass through the model.

%         Args:
%             x: Input tensor of shape (batch_size, seq_len).
%             lengths: Real sequence lengths before padding (batch_size,).
%                     If None, assumes all sequences are used fully.

%         Returns:
%             Logits tensor of shape (batch_size, num_classes).
%         """
%         embedded = self.embedding(x)

%         if lengths is not None:
%             # Pack sequences to ignore padding tokens
%             packed = nn.utils.rnn.pack_padded_sequence(
%                 embedded, lengths.cpu(), batch_first=True, enforce_sorted=False
%             )
%             _, (hidden, _) = self.lstm(packed)
%             # hidden shape: (num_layers * num_directions, batch, hidden_dim)
%         else:
%             # No packing - process all positions
%             _, (hidden, _) = self.lstm(embedded)

%         # Extract final hidden state
%         if self.lstm.bidirectional:
%             # Concatenate final forward and backward hidden states
%             # hidden[-2] is the last forward layer, hidden[-1] is the last backward layer
%             last_hidden = torch.cat([hidden[-2], hidden[-1]], dim=1)
%         else:
%             # Just take the last layer's hidden state
%             last_hidden = hidden[-1]

%         # Apply dropout and classification layer
%         return self.fc(self.dropout(last_hidden))

